\documentclass[11pt]{article}

\usepackage[a4paper,margin=27mm]{geometry}
\usepackage{amsmath,amssymb,amsthm,mathtools}
\usepackage{booktabs,longtable,array}
\usepackage[hidelinks]{hyperref}
\usepackage{microtype}
\usepackage{enumitem}

\newtheorem{definition}{Definition}[section]
\newtheorem{proposition}[definition]{Proposition}
\newtheorem{theorem}[definition]{Theorem}
\newtheorem{corollary}[definition]{Corollary}
\newtheorem{example}[definition]{Example}
\newtheorem{remark}[definition]{Remark}

\newcommand{\W}{\mathcal W}
\newcommand{\Pcal}{\mathcal P}
\newcommand{\Acal}{\mathcal A}
\newcommand{\Bcal}{\mathcal B}
\newcommand{\Y}{\mathcal Y}
\newcommand{\Fin}{\operatorname{Fin}}
\newcommand{\Cfin}{\mathfrak C_{\mathrm{fin}}}
\newcommand{\Dfin}{\mathfrak D_{\mathrm{fin}}}

\hypersetup{
  pdftitle={WR-II: Finite Consistency and Global World Realizability},
  pdfauthor={A. A. Malachevsky},
  pdfsubject={World Realizability and Epistemic Horizons (WREH), WR-II},
  pdfkeywords={world realizability, finite consistency, inverse limits, compactness, response profiles, global realization, epistemic horizons}
}

\title{\textbf{WR-II Preprint Draft v0.3}\\[4pt]
\large Finite Consistency and Global World Realizability}
\author{A. A. Malachevsky\\
ORCID: 0009-0008-6009-3196\\
World Realizability \& Epistemic Horizons (WREH)\\
Research Programme \& Community}
\date{7 October 2026}

\begin{document}
\maketitle

\begin{abstract}
WR-II studies the local-to-global question left open by WR-I: whether a response
profile whose every finite restriction is realized by some admissible global
realization must itself be realized by one admissible global realization.
Starting from the WR-I world-response structure, we construct an inverse system
of finite response images and define its inverse limit as the finite-consistency
space. The full WR-I response quotient embeds canonically into this space; the
complement of its image is the global realizability defect. We show that this
defect is impossible for a finite protocol family, exhibit an explicit
infinite-protocol counterexample based on finite-support binary sequences, and
identify the finite-consistency space with a finite-projection closure of the
full response image. This yields a closed-image criterion and, in particular,
a compactness criterion for global realizability. We then introduce
protocol-normalized residual radii for noisy response profiles and prove a
finite-to-global residual identity on compact world spaces, together with a
compact-core extension that prevents escape to infinity. Noncompact examples
separate positive stability gaps from zero-distance nonattainment. For
countable protocol families we also give nested-fibre existence and uniqueness
criteria. Finally, we distinguish full finite consistency from
context-restricted local consistency, where local-to-global failure can already
occur for finitely many protocols. The framework is model-relative and does
not identify mathematical completion with physical existence.
\end{abstract}

\section{Upstream interface and claim ceiling}

WR-II begins from the world-response structure fixed in WR-I
\cite{MalachevskyWRI}:
\[
\mathbf W=(\W,\Pcal,\{R_P\}_{P\in\Pcal}),
\]
where $\W$ is a declared class of admissible global realizations, $\Pcal$ is a
declared family of admissible protocols, and
\[
R_P:\W\to Y_P
\]
is the response map associated with protocol $P$.

WR-I studied identifiability relative to a fixed protocol family. WR-II asks a
different question: whether a family of locally or finitely realizable response
constraints admits one common global realization.

The claim ceiling is strict. WR-II does not claim that:
\begin{itemize}
\item mathematical finite consistency implies physical existence;
\item a profile not realized inside the declared class $\W$ is physically
impossible in an absolute sense;
\item a formal inverse-limit point is itself a physical world;
\item global completion is unique without a separate uniqueness argument;
\item refinement order is physical time, dynamics, or renormalization flow;
\item failure of global realizability means that the physical Universe fails
to exist.
\end{itemize}

\section{Full response profiles and finite response images}

Define the full response product
\[
\Y:=\prod_{P\in\Pcal}Y_P
\]
and the full response map
\[
R:\W\to\Y,
\qquad
R(W):=(R_P(W))_{P\in\Pcal}.
\]

For a finite bundle $\Acal\in\Fin(\Pcal)$, let
\[
Y_{\Acal}:=\prod_{P\in\Acal}Y_P,
\qquad
R_{\Acal}:=\pi_{\Acal}\circ R,
\]
where $\pi_{\Acal}:\Y\to Y_{\Acal}$ is the coordinate projection.

\begin{definition}[Finite realizable image]
For $\Acal\in\Fin(\Pcal)$ define
\[
E_{\Acal}:=R_{\Acal}(\W)\subseteq Y_{\Acal}.
\]
\end{definition}

If $\Acal\subseteq\Bcal$, coordinate restriction gives
\[
\rho_{\Bcal\Acal}:E_{\Bcal}\to E_{\Acal}.
\]

\begin{proposition}[Finite response inverse system]\label{prop:inverse-system}
For every $\Acal\subseteq\Bcal$, the restriction map
$\rho_{\Bcal\Acal}$ is surjective. Moreover,
\[
\rho_{\Acal\Acal}=\mathrm{id}_{E_{\Acal}},
\qquad
\rho_{\mathcal C\Acal}
=
\rho_{\Bcal\Acal}\circ\rho_{\mathcal C\Bcal}
\]
whenever $\Acal\subseteq\Bcal\subseteq\mathcal C$.
Hence $(E_{\Acal},\rho_{\Bcal\Acal})$ is an inverse system indexed by
$\Fin(\Pcal)$.
\end{proposition}

\begin{proof}
Take $e_{\Acal}\in E_{\Acal}$. By definition there is $W\in\W$ with
$R_{\Acal}(W)=e_{\Acal}$. Then $R_{\Bcal}(W)\in E_{\Bcal}$ restricts to
$e_{\Acal}$. The identity and composition laws are coordinate restriction.
\end{proof}

\section{Finite-consistency space}

\begin{definition}[Finite-consistency space]
The \emph{finite-consistency space} of $\mathbf W$ is
\[
\Cfin(\mathbf W)
:=
\varprojlim_{\Acal\in\Fin(\Pcal)}E_{\Acal}.
\]
Its points are compatible families
\[
e=(e_{\Acal})_{\Acal\in\Fin(\Pcal)}
\]
with
\[
\rho_{\Bcal\Acal}(e_{\Bcal})=e_{\Acal}
\qquad
(\Acal\subseteq\Bcal).
\]
\end{definition}

There is an equivalent profile description.

\begin{proposition}[Profile representation]\label{prop:profile-representation}
There is a canonical bijection
\[
\Cfin(\mathbf W)
\cong
\left\{
y\in\Y:
\pi_{\Acal}(y)\in E_{\Acal}
\text{ for every finite }\Acal\subseteq\Pcal
\right\}.
\]
Thus a point of $\Cfin$ is exactly a full response profile whose every finite
restriction is realized by at least one admissible global realization.
\end{proposition}

\begin{proof}
A compatible inverse-limit family determines its singleton coordinates and
therefore a unique point $y\in\Y$. Compatibility forces every finite component
to equal the corresponding finite coordinate tuple of $y$. Conversely, any
such $y$ defines a compatible family by finite restriction.
\end{proof}

\section{Canonical embedding of the WR-I quotient}

Let
\[
\mathcal Q_{\Pcal}:=\W/\!\sim_{\Pcal}
\]
be the full WR-I response quotient.

\begin{proposition}[Canonical global-realization embedding]
There is a canonical injection
\[
J:\mathcal Q_{\Pcal}\hookrightarrow\Cfin(\mathbf W),
\qquad
J([W])=(R_{\Acal}(W))_{\Acal\in\Fin(\Pcal)}.
\]
Under Proposition~\ref{prop:profile-representation}, the image of $J$ is
exactly the full response image $R(\W)\subseteq\Y$.
\end{proposition}

\begin{proof}
If $[W]=[W']$ in $\mathcal Q_{\Pcal}$, all protocol responses agree, hence all
finite responses agree, so $J$ is well defined. If $J([W])=J([W'])$, singleton
components agree for every $P\in\Pcal$, so $W\sim_{\Pcal}W'$ and $J$ is
injective. The image statement is immediate.
\end{proof}

\begin{definition}[Global realizability defect]
The \emph{global realizability defect} is
\[
\Dfin(\mathbf W)
:=
\Cfin(\mathbf W)\setminus J(\mathcal Q_{\Pcal}).
\]
\end{definition}

This definition is model-relative: $\Dfin$ measures failure of realization
inside the declared class $\W$.

\section{Three-way realizability classification}

For any proposed full response profile $y\in\Y$, exactly one of the following
holds.

\begin{proposition}[Profile trichotomy]
Each $y\in\Y$ is in exactly one of three states:
\begin{enumerate}[label=(\roman*)]
\item \textbf{finite obstruction:} there exists finite $\Acal\subseteq\Pcal$
with $\pi_{\Acal}(y)\notin E_{\Acal}$;
\item \textbf{global completion defect:}
$y\in\Cfin(\mathbf W)\setminus R(\W)$;
\item \textbf{global realization:} $y\in R(\W)$.
\end{enumerate}
\end{proposition}

\begin{proof}
By Proposition~\ref{prop:profile-representation}, absence of a finite
obstruction is equivalent to membership in $\Cfin$. Since
$R(\W)\subseteq\Cfin$, the remaining two cases partition $\Cfin$.
\end{proof}

\section{Why the protocol family must be infinite for a genuine full-finite defect}

\begin{proposition}[Finite-family collapse]\label{prop:finite-collapse}
If $\Pcal$ is finite, then
\[
\Cfin(\mathbf W)=R(\W)
\]
and therefore
\[
\Dfin(\mathbf W)=\varnothing.
\]
\end{proposition}

\begin{proof}
Since $\Pcal\in\Fin(\Pcal)$, finite consistency includes the requirement
\[
\pi_{\Pcal}(y)=y\in E_{\Pcal}=R(\W).
\]
\end{proof}

\begin{remark}
This elementary observation is structurally important. A genuine distinction
between \emph{all finite bundles} and a global response profile requires an
infinite protocol family. Finite contextuality scenarios avoid this collapse
because only specified jointly admissible contexts, not the full finite power
set, are required to be locally realizable.
\end{remark}

\section{Explicit failure of finite-to-global realizability}

\begin{example}[Finite-support binary worlds]\label{ex:finite-support}
Let
\[
\Pcal=\mathbb N,
\qquad
Y_n=\{0,1\},
\]
and let
\[
\W
=
\bigoplus_{n\in\mathbb N}\mathbb Z_2
=
\left\{
x\in\{0,1\}^{\mathbb N}:
x_n=1\text{ for only finitely many }n
\right\}.
\]
Let protocol $P_n$ read the $n$th coordinate:
\[
R_n(x)=x_n.
\]

For every finite $\Acal\subset\mathbb N$,
\[
E_{\Acal}=\{0,1\}^{\Acal},
\]
because any finite binary pattern can be extended by zeros to an element of
$\W$. Hence
\[
\Cfin(\mathbf W)=\{0,1\}^{\mathbb N}.
\]
But
\[
R(\W)=\W
\]
contains only finite-support profiles. Therefore
\[
\Dfin(\mathbf W)
=
\{0,1\}^{\mathbb N}\setminus
\bigoplus_{n\in\mathbb N}\mathbb Z_2.
\]

In particular, the profile
\[
y=(1,1,1,\ldots)
\]
is realized on every finite protocol bundle but by no world in $\W$.
\end{example}

This example shows that finite realizability alone does not imply a common
global realization.

\section{Finite-projection closure}

The previous construction has a simple closure formulation.

\begin{definition}[Finite-projection closure]
For $S\subseteq\Y$, define
\[
\operatorname{Cl}_{\mathrm{fin}}(S)
:=
\bigcap_{\Acal\in\Fin(\Pcal)}
\pi_{\Acal}^{-1}\!\left(\pi_{\Acal}(S)\right).
\]
\end{definition}

\begin{proposition}[Closure-operator properties]
The map $S\mapsto\operatorname{Cl}_{\mathrm{fin}}(S)$ is extensive,
monotone, and idempotent:
\[
S\subseteq\operatorname{Cl}_{\mathrm{fin}}(S),
\]
\[
S\subseteq T
\Longrightarrow
\operatorname{Cl}_{\mathrm{fin}}(S)
\subseteq
\operatorname{Cl}_{\mathrm{fin}}(T),
\]
and
\[
\operatorname{Cl}_{\mathrm{fin}}
\!\left(\operatorname{Cl}_{\mathrm{fin}}(S)\right)
=
\operatorname{Cl}_{\mathrm{fin}}(S).
\]
\end{proposition}

\begin{proof}
Extensivity and monotonicity are immediate. If
$C=\operatorname{Cl}_{\mathrm{fin}}(S)$, then for every finite $\Acal$,
\[
\pi_{\Acal}(C)=\pi_{\Acal}(S).
\]
The inclusion $\pi_{\Acal}(S)\subseteq\pi_{\Acal}(C)$ follows from
$S\subseteq C$, while the reverse inclusion is part of the definition of $C$.
Applying the definition again therefore changes nothing.
\end{proof}

\begin{proposition}[Finite consistency as completion]\label{prop:completion}
For the full response image $S=R(\W)$,
\[
\Cfin(\mathbf W)
=
\operatorname{Cl}_{\mathrm{fin}}(R(\W)).
\]
Consequently,
\[
\Dfin(\mathbf W)=\varnothing
\quad\Longleftrightarrow\quad
R(\W)
=
\operatorname{Cl}_{\mathrm{fin}}(R(\W)).
\]
\end{proposition}

\begin{remark}
No novelty claim is attached to the abstract idea of closure by finite
projections. WR-II uses this operator because it exposes the precise gap
between finite response realizability and realization by a single world.
\end{remark}

\section{Topological criteria for global realizability}

Give each $Y_P$ a topology and $\Y$ the product topology.

\begin{proposition}[Finite-projection completion lies in topological closure]
For every $S\subseteq\Y$,
\[
\operatorname{Cl}_{\mathrm{fin}}(S)
\subseteq
\overline S.
\]
\end{proposition}

\begin{proof}
Let $y\in\operatorname{Cl}_{\mathrm{fin}}(S)$. Every basic neighborhood of
$y$ restricts only finitely many coordinates, say those in $\Acal$. Since
$\pi_{\Acal}(y)\in\pi_{\Acal}(S)$, some $s\in S$ agrees with $y$ on all
coordinates in $\Acal$, hence lies in that basic neighborhood. Therefore every
neighborhood of $y$ meets $S$.
\end{proof}

\begin{theorem}[Closed-image criterion]\label{thm:closed-image}
If the full response image $R(\W)$ is closed in the product topology on $\Y$,
then
\[
\Cfin(\mathbf W)=R(\W)
\]
and the global realizability defect is empty.
\end{theorem}

\begin{proof}
By Proposition~\ref{prop:completion},
\[
\Cfin=\operatorname{Cl}_{\mathrm{fin}}(R(\W))
\subseteq
\overline{R(\W)}
=
R(\W).
\]
The reverse inclusion always holds.
\end{proof}

\begin{corollary}[Compact-domain global realizability]\label{cor:compact}
Assume:
\begin{enumerate}[label=(\roman*)]
\item $\W$ is compact;
\item each response space $Y_P$ is Hausdorff;
\item each $R_P:\W\to Y_P$ is continuous.
\end{enumerate}
Then
\[
\Dfin(\mathbf W)=\varnothing.
\]
\end{corollary}

\begin{proof}
The full map
\[
R:\W\to\prod_{P\in\Pcal}Y_P
\]
is continuous. The product is Hausdorff, so the compact set $R(\W)$ is closed.
Apply Theorem~\ref{thm:closed-image}.
\end{proof}

\begin{remark}
This is a sufficient condition, not a claim that physically admissible world
spaces are compact. Noncompact models may still be globally realizable, while
Example~\ref{ex:finite-support} shows that noncompactness can permit a genuine
completion defect.
\end{remark}


\section{Noisy finite consistency and quantitative global realizability}

Exact equality is too strong for experimental data. Assume from now on that
each response space $Y_P$ is equipped with a metric $d_P$, and choose a
strictly positive protocol scale
\[
\sigma_P>0.
\]
The scale may represent a declared resolution, tolerance, or normalization
needed to compare heterogeneous protocols. It is part of the model and is not
canonically determined by the abstract world-response structure. For the
topological results below, $\W$ is additionally equipped with a topology.

Fix a target profile
\[
y=(y_P)_{P\in\Pcal}\in\Y.
\]

\begin{definition}[Normalized protocol residual]
Define
\[
r_P(W;y)
:=
\frac{d_P(R_P(W),y_P)}{\sigma_P}.
\]
For finite $\Acal\subseteq\Pcal$, define
\[
r_{\Acal}(W;y)
:=
\max_{P\in\Acal}r_P(W;y).
\]
\end{definition}


\begin{definition}[Weighted uniform response discrepancy]
For profiles $y,z\in\Y$, define the extended discrepancy
\[
d_{\sigma,\infty}(y,z)
:=
\sup_{P\in\Pcal}
\frac{d_P(y_P,z_P)}{\sigma_P}
\in[0,\infty].
\]
For finite $\Acal$, define
\[
d_{\sigma,\Acal}
(\pi_{\Acal}y,\pi_{\Acal}z)
:=
\max_{P\in\Acal}
\frac{d_P(y_P,z_P)}{\sigma_P}.
\]
\end{definition}

\begin{definition}[Finite and global residual radii]
For finite $\Acal$, let
\[
\delta_{\Acal}(y)
:=
\inf_{W\in\W} r_{\Acal}(W;y).
\]
Define
\[
\delta_{\mathrm{fin}}(y)
:=
\sup_{\Acal\in\Fin(\Pcal)}
\delta_{\Acal}(y)
\]
and
\[
\delta_{\mathrm{glob}}(y)
:=
\inf_{W\in\W}
\sup_{P\in\Pcal}r_P(W;y).
\]
These are extended nonnegative numbers.
\end{definition}

\begin{proposition}[Distance-to-response-image representation]
For every target profile $y$,
\[
\delta_{\mathrm{glob}}(y)
=
\inf_{z\in R(\W)}d_{\sigma,\infty}(y,z).
\]
For every finite $\Acal$,
\[
\delta_{\Acal}(y)
=
\inf_{z_{\Acal}\in E_{\Acal}}
d_{\sigma,\Acal}
(\pi_{\Acal}y,z_{\Acal}),
\]
and therefore
\[
\delta_{\mathrm{fin}}(y)
=
\sup_{\Acal\in\Fin(\Pcal)}
\operatorname{dist}_{\sigma,\Acal}
\bigl(\pi_{\Acal}y,E_{\Acal}\bigr).
\]
\end{proposition}

\begin{proof}
Each formula is the corresponding definition rewritten through the image of
the response map.
\end{proof}

The finite radius asks how well every finite response bundle can be fitted,
allowing the fitting world to depend on the bundle. The global radius asks for
one world that fits all protocols simultaneously.

\begin{proposition}[Finite-to-global residual inequality]
For every target profile $y$,
\[
\delta_{\mathrm{fin}}(y)
\le
\delta_{\mathrm{glob}}(y).
\]
\end{proposition}

\begin{proof}
For every $W\in\W$ and every finite $\Acal$,
\[
r_{\Acal}(W;y)
\le
\sup_{P\in\Pcal}r_P(W;y).
\]
Taking the infimum over $W$ and then the supremum over finite $\Acal$ gives the
claim.
\end{proof}

\begin{definition}[Finite-to-global stability gap]
When both radii are finite, define
\[
\Gamma(y)
:=
\delta_{\mathrm{glob}}(y)
-
\delta_{\mathrm{fin}}(y)
\ge0.
\]
A positive value measures a quantitative failure of finite fits to globalize
at the same normalized tolerance.
\end{definition}

The numerical values of the radii depend on the declared protocol
normalization. Nevertheless, boundedly equivalent calibrations preserve their
zero-versus-positive character.

\begin{proposition}[Uniform calibration stability]
Let $(\sigma_P)$ and $(\tau_P)$ be two positive protocol scales satisfying
\[
0<c\le
\frac{\tau_P}{\sigma_P}
\le C<\infty
\qquad
\text{for every }P\in\Pcal.
\]
Let $\delta^{\sigma}_{\mathrm{fin}},
\delta^{\sigma}_{\mathrm{glob}}$ and
$\delta^{\tau}_{\mathrm{fin}},
\delta^{\tau}_{\mathrm{glob}}$
denote the radii computed with the two calibrations. Then
\[
\frac1C\,
\delta^{\sigma}_{\mathrm{fin}}(y)
\le
\delta^{\tau}_{\mathrm{fin}}(y)
\le
\frac1c\,
\delta^{\sigma}_{\mathrm{fin}}(y),
\]
and
\[
\frac1C\,
\delta^{\sigma}_{\mathrm{glob}}(y)
\le
\delta^{\tau}_{\mathrm{glob}}(y)
\le
\frac1c\,
\delta^{\sigma}_{\mathrm{glob}}(y).
\]
In particular, whether either radius is zero is invariant under uniformly
equivalent changes of scale.
\end{proposition}

\begin{proof}
The scale bounds imply, for every protocol,
\[
\frac1C r_P^\sigma(W;y)
\le
r_P^\tau(W;y)
\le
\frac1c r_P^\sigma(W;y).
\]
The same inequalities survive finite maxima, arbitrary suprema, and infima over
$W$.
\end{proof}

\begin{remark}
The hypothesis of uniform equivalence is essential. The scale choice
$\sigma_n=n$ in the zero-margin nonattainment example is not uniformly
equivalent to the unit calibration; it deliberately changes the asymptotic
importance of high-index protocols.
\end{remark}

\begin{theorem}[Compact finite-to-global residual identity]
\label{thm:robust-compact}
Assume $\W$ is compact and, for every $P\in\Pcal$, the residual function
\[
W\longmapsto r_P(W;y)
\]
is lower semicontinuous. If $\delta_{\mathrm{fin}}(y)<\infty$, then
\[
\delta_{\mathrm{glob}}(y)
=
\delta_{\mathrm{fin}}(y),
\]
and the global infimum is attained by at least one $W_\ast\in\W$.
\end{theorem}

\begin{proof}
Write
\[
\alpha:=\delta_{\mathrm{fin}}(y).
\]
For finite $\Acal$, the function
\[
r_{\Acal}(\cdot;y)
=
\max_{P\in\Acal}r_P(\cdot;y)
\]
is lower semicontinuous. Compactness implies that it attains its minimum
$\delta_{\Acal}(y)$.

Define the closed nonempty set
\[
K_{\Acal}
:=
\{W\in\W:r_{\Acal}(W;y)\le\alpha\}.
\]
The family $(K_{\Acal})_{\Acal\in\Fin(\Pcal)}$ has the finite intersection
property: for finite bundles $\Acal_1,\ldots,\Acal_m$, let
\[
\Bcal:=\Acal_1\cup\cdots\cup\Acal_m.
\]
Then $K_{\Bcal}$ is nonempty and
\[
K_{\Bcal}\subseteq
K_{\Acal_1}\cap\cdots\cap K_{\Acal_m}.
\]
By compactness,
\[
\bigcap_{\Acal\in\Fin(\Pcal)}K_{\Acal}\neq\varnothing.
\]
Choose $W_\ast$ in the intersection. In particular, for every singleton
$\{P\}$,
\[
r_P(W_\ast;y)\le\alpha.
\]
Hence
\[
\sup_{P\in\Pcal}r_P(W_\ast;y)\le\alpha,
\]
so
\[
\delta_{\mathrm{glob}}(y)\le\alpha.
\]
The reverse inequality is the preceding proposition.
\end{proof}

\begin{corollary}[Tolerance globalization]
Under the hypotheses of Theorem~\ref{thm:robust-compact}, suppose that for
every finite $\Acal\subseteq\Pcal$ there exists $W_{\Acal}\in\W$ such that
\[
d_P(R_P(W_{\Acal}),y_P)\le\sigma_P
\qquad
\text{for all }P\in\Acal.
\]
Then there exists one $W_\ast\in\W$ such that
\[
d_P(R_P(W_\ast),y_P)\le\sigma_P
\qquad
\text{for every }P\in\Pcal.
\]
\end{corollary}

\begin{proof}
The hypothesis gives $\delta_{\Acal}(y)\le1$ for every finite $\Acal$, so
$\delta_{\mathrm{fin}}(y)\le1$. Apply
Theorem~\ref{thm:robust-compact}.
\end{proof}

\begin{corollary}[Vanishing-noise exactification]
Under the hypotheses of Theorem~\ref{thm:robust-compact}, if
\[
\delta_{\mathrm{fin}}(y)=0,
\]
then there exists $W_\ast\in\W$ with
\[
R_P(W_\ast)=y_P
\qquad
\text{for every }P\in\Pcal.
\]
\end{corollary}

\begin{proof}
The theorem gives an attained global radius equal to zero. Since every metric
is definite, every protocol residual vanishes.
\end{proof}

Global compactness is stronger than necessary. A finite bundle can itself
provide the compactness needed to prevent approximate witnesses from escaping
to infinity.

\begin{theorem}[Compact-core residual identity]
\label{thm:compact-core}
Assume every residual $r_P(\cdot;y)$ is lower semicontinuous and
\[
\alpha:=\delta_{\mathrm{fin}}(y)<\infty.
\]
Suppose there exist a finite bundle $\Acal_0$ and a number $c>\alpha$ such
that the sublevel set
\[
K
:=
\{W\in\W:r_{\Acal_0}(W;y)\le c\}
\]
is compact. Then
\[
\delta_{\mathrm{glob}}(y)=\alpha,
\]
and the global infimum is attained.
\end{theorem}

\begin{proof}
Choose $N$ such that
\[
\alpha+\frac1N<c.
\]
For finite $\Acal$ and $n\ge N$, define
\[
K_{\Acal,n}
:=
\left\{
W\in K:
r_{\Acal}(W;y)\le\alpha+\frac1n
\right\}.
\]
These sets are closed in $K$.

They are nonempty. Indeed,
\[
\delta_{\Acal\cup\Acal_0}(y)\le\alpha,
\]
so for every $n\ge N$ there exists $W$ with
\[
r_{\Acal\cup\Acal_0}(W;y)
\le
\alpha+\frac1n.
\]
Such a point lies in $K$ because
\[
r_{\Acal_0}(W;y)
\le
\alpha+\frac1n
<c.
\]

The family $\{K_{\Acal,n}\}$ has the finite intersection property. For a
finite collection $(\Acal_i,n_i)$, let
\[
\Bcal:=\Acal_0\cup\bigcup_i\Acal_i,
\qquad
m:=\max_i n_i.
\]
Choose $W$ with
\[
r_{\Bcal}(W;y)\le\alpha+\frac1m.
\]
Then $W$ belongs to every selected $K_{\Acal_i,n_i}$.

Compactness of $K$ gives a point $W_\ast$ in the intersection of all
$K_{\Acal,n}$ for finite $\Acal$ and $n\ge N$. Letting $n\to\infty$ yields
\[
r_{\Acal}(W_\ast;y)\le\alpha
\]
for every finite $\Acal$. Hence
\[
\sup_{P\in\Pcal}r_P(W_\ast;y)\le\alpha.
\]
Combine with the general inequality.
\end{proof}

\begin{remark}
The compact-core theorem is a coercivity principle. The entire world class
need not be compact: it is enough that one finite response bundle traps all
near-optimal global candidates inside a compact sublevel set.
\end{remark}

\begin{example}[Positive stability gap without compactness]
Return to the finite-support binary world class of
Example~\ref{ex:finite-support}. Give every $Y_n=\{0,1\}$ the discrete metric,
set $\sigma_n=1$, and use the all-ones target
\[
y=(1,1,\ldots).
\]
Every finite bundle is fitted exactly, so
\[
\delta_{\mathrm{fin}}(y)=0.
\]
Every finite-support world has at least one zero coordinate, hence
\[
\sup_n d_n(R_n(W),1)=1.
\]
Therefore
\[
\delta_{\mathrm{glob}}(y)=1,
\qquad
\Gamma(y)=1.
\]
\end{example}

\begin{example}[Zero-margin nonattainment and scale dependence]
Use the same finite-support world class and the same all-ones target, but choose
protocol scales
\[
\sigma_n=n.
\]
Again
\[
\delta_{\mathrm{fin}}(y)=0.
\]
For the world having ones in the first $N$ coordinates and zeros afterwards,
\[
\sup_n r_n(W;y)
=
\frac{1}{N+1}.
\]
Hence
\[
\delta_{\mathrm{glob}}(y)=0,
\]
but no world realizes $y$ exactly.

Thus a zero global residual radius does not imply exact realization when the
infimum is not attained. Moreover, the quantitative margin depends on the
declared protocol metrics and scales. Any physical use of $\Gamma$ therefore
requires an independently justified calibration of those choices.
\end{example}

\begin{remark}[Not an approximate inverse system]
The inverse system of finite response images remains exactly commutative in
this section. Noise is introduced in the target response profile and in its
metric residuals, not in the bonding maps. This differs from the classical
theory of approximate inverse systems, where commutativity of the bonding
maps is itself relaxed in a controlled way
\cite{MardesicRubin1989,MardesicSegal1989}.
\end{remark}

\begin{remark}[Relation to robust contextuality]
For context-restricted probability models, quantitative distances and
robustness measures relative to the noncontextual set are already well
developed, including the contextual fraction
\cite{AbramskyBarbosaMansfield2017}. WR-II does not claim novelty for
quantitative contextuality. The residual radii above are defined for an
arbitrary world-response structure and address a different question: whether
finite protocol fits at a declared tolerance can be globalized by one world.
\end{remark}


\section{Nested-fibre criteria for countable protocol families}

Assume now that $\Pcal=\{P_1,P_2,\ldots\}$ is countable and define
\[
\Acal_n:=\{P_1,\ldots,P_n\}.
\]
For a finitely consistent profile $y\in\Cfin$, define
\[
F_n(y)
:=
\{W\in\W:R_{\Acal_n}(W)=\pi_{\Acal_n}(y)\}.
\]
Then
\[
F_1(y)\supseteq F_2(y)\supseteq\cdots
\]
and every $F_n(y)$ is nonempty.

\begin{proposition}[Compact nested-fibre existence]
Assume $\W$ is compact and every fibre $F_n(y)$ is closed in $\W$. Then
\[
\bigcap_{n=1}^{\infty}F_n(y)\neq\varnothing.
\]
Every point in this intersection globally realizes $y$.
\end{proposition}

\begin{proof}
The sets $F_n(y)$ are nonempty, closed, and nested. Compactness gives a
nonempty intersection. Agreement on every $\Acal_n$ implies agreement on every
protocol.
\end{proof}

\begin{corollary}[Continuous Hausdorff-response version]
If $\W$ is compact, every $Y_P$ is Hausdorff, and every response
$R_P:\W\to Y_P$ is continuous, then every finitely consistent profile for a
countable protocol family has at least one global realization.
\end{corollary}

\begin{proof}
For each $n$, the target point $\pi_{\Acal_n}(y)$ is closed in the finite
Hausdorff product $Y_{\Acal_n}$. Hence
\[
F_n(y)=R_{\Acal_n}^{-1}\!\left(\{\pi_{\Acal_n}(y)\}\right)
\]
is closed. Apply the preceding proposition.
\end{proof}

\begin{theorem}[Metric shrinking-fibre uniqueness]\label{thm:shrinking}
Let $(\W,d)$ be complete. Suppose each $F_n(y)$ is nonempty and closed and
\[
\operatorname{diam}F_n(y)\longrightarrow0.
\]
Then
\[
\bigcap_{n=1}^{\infty}F_n(y)
=
\{W_y\}
\]
for a unique $W_y\in\W$, and $R(W_y)=y$.
\end{theorem}

\begin{proof}
Choose $W_n\in F_n(y)$. If $m\ge n$, then both $W_m$ and $W_n$ lie in
$F_n(y)$, hence
\[
d(W_m,W_n)\le\operatorname{diam}F_n(y)\to0.
\]
Thus $(W_n)$ is Cauchy and converges to some $W_y\in\W$. For fixed $n$, all
$W_m$ with $m\ge n$ lie in the closed set $F_n(y)$, so $W_y\in F_n(y)$.
Therefore $W_y$ belongs to every $F_n(y)$. If $W'$ also belongs to every
$F_n(y)$, then
\[
d(W_y,W')\le\operatorname{diam}F_n(y)
\]
for all $n$, hence $W'=W_y$.
\end{proof}

\section{Context-restricted local consistency}

Full finite consistency uses every finite bundle. Physical theories may instead
specify only certain jointly admissible contexts.

\begin{definition}[Context family]
A \emph{context family} is a collection
\[
\mathfrak C\subseteq\Fin(\Pcal)
\]
covering $\Pcal$ and closed under taking subsets. For each
$C\in\mathfrak C$, define
\[
E_C:=R_C(\W).
\]
\end{definition}

The inclusion order and coordinate restrictions form a diagram of local
response images even when $\mathfrak C$ is not directed.

\begin{definition}[Context-consistency space]
Define
\[
\mathfrak C_{\mathrm{ctx}}(\mathbf W,\mathfrak C)
:=
\varprojlim_{C\in\mathfrak C}E_C,
\]
that is, the set of families $(e_C)_{C\in\mathfrak C}$ satisfying
\[
e_C\in E_C
\]
and agreeing under every restriction $D\subseteq C$ in $\mathfrak C$.
\end{definition}

Because $\mathfrak C$ covers $\Pcal$, there is a canonical map
\[
J_{\mathfrak C}:\mathcal Q_{\Pcal}
\longrightarrow
\mathfrak C_{\mathrm{ctx}}(\mathbf W,\mathfrak C)
\]
sending a world class to all of its contextual responses.

\begin{proposition}[Context embedding]
If $\mathfrak C$ covers $\Pcal$ and is closed under subsets, then
$J_{\mathfrak C}$ is injective.
\end{proposition}

\begin{proof}
If two world classes have the same contextual family, then for every
$P\in\Pcal$ their responses agree on the singleton context $\{P\}$.
Hence they are equivalent under the full protocol family.
\end{proof}

\begin{definition}[Contextual global-realizability defect]
The context-relative defect is
\[
\mathfrak D_{\mathrm{ctx}}(\mathbf W,\mathfrak C)
:=
\mathfrak C_{\mathrm{ctx}}(\mathbf W,\mathfrak C)
\setminus
J_{\mathfrak C}(\mathcal Q_{\Pcal}).
\]
\end{definition}

\begin{example}[Finite contextual completion defect]
Let
\[
\Pcal=\{1,2,3\},
\qquad
\W=\{000,011,101,110\},
\]
the even-parity binary triples. Let the maximal contexts be
\[
\{1,2\},\qquad \{2,3\},\qquad \{1,3\},
\]
together with all their subsets.

Every binary pair occurs as the projection of some $W\in\W$. Hence the local
pair responses
\[
(1,1)\text{ on }\{1,2\},\quad
(1,1)\text{ on }\{2,3\},\quad
(1,1)\text{ on }\{1,3\}
\]
are individually realizable and agree on all overlaps. They define the global
coordinate profile $111$, but
\[
111\notin\W.
\]
Thus context-local consistency fails to extend to a world realization even
though $\Pcal$ is finite.
\end{example}

\begin{remark}
This does not contradict Proposition~\ref{prop:finite-collapse}: the full
bundle $\{1,2,3\}$ is not a member of the chosen context family. This is the
precise interface with contextuality and other local-to-global obstruction
theories.
\end{remark}

\section{Prior-art boundary}

WR-II sits next to several mature finite-to-global theories.

\begin{longtable}{>{\raggedright\arraybackslash}p{0.22\textwidth}
                  >{\raggedright\arraybackslash}p{0.34\textwidth}
                  >{\raggedright\arraybackslash}p{0.34\textwidth}}
\toprule
\textbf{Prior line} & \textbf{Established result} & \textbf{WR-II boundary}\\
\midrule
\endfirsthead
\toprule
\textbf{Prior line} & \textbf{Established result} & \textbf{WR-II boundary}\\
\midrule
\endhead

First-order compactness &
If every finite subset of a first-order theory is satisfiable, the whole theory
is satisfiable. &
WR-II does not rebrand logical compactness. Its constraints are response images
of a declared world class, with topology and admissibility supplied by the
application. \\[5pt]

Inverse/projective limits &
Classical inverse-limit theory supplies existence results under compactness,
surjectivity, Mittag--Leffler, or related hypotheses. &
WR-II uses the inverse limit of finite response images as a diagnostic space and
compares it with the image of the WR-I full response quotient. \\[5pt]

Kolmogorov extension &
Compatible finite-dimensional probability laws admit a global stochastic
process under the standard hypotheses. &
WR-II response profiles need not be probability measures. Kolmogorov extension
is a domain-specific positive archetype, not the WR-II theorem. \\[5pt]

Vorob'ev (1962) &
Characterizes extendability of consistent families of measures through the
structure of the underlying complex \cite{Vorobev1962}. &
Direct predecessor for context-restricted extension questions. WR-II keeps
world realizations distinct from probability measures. \\[5pt]

Abramsky--Brandenburger (2011) &
Contextuality and non-locality correspond to obstructions to global sections
over measurement covers \cite{AbramskyBrandenburger2011}. &
Direct predecessor for the context-family layer. WR-II does not identify a
global section with a physical world without an explicit bridge. \\[5pt]

Atserias--Kolaitis (2023) &
Acyclicity characterizes when local consistency of positive-semiring relations
forces global consistency, generalizing probability and database results
\cite{AtseriasKolaitis2023}. &
Shows that local-to-global consistency has a broad established theory. WR-II's
claim is architectural and model-relative, not theorem-priority over that
literature. \\[5pt]

Semi-infinite programming &
Optimization with finitely many decision variables and infinitely many
constraints is a mature field, with compactness, discretization, duality, and
exchange methods surveyed by Hettich and Kortanek
\cite{HettichKortanek1993}. &
The noisy WR-II global residual is a response-space fitting functional and can
be viewed as a semi-infinite min--max problem in suitable applications. WR-II
does not claim priority for infinite-constraint optimization. \\
\bottomrule
\end{longtable}

\section{Novelty claim permitted at v0.1}

The defensible claim at this stage is architectural:

\begin{quote}
WR-II attaches to the WR-I world-response structure a canonical finite-response
inverse system, identifies its inverse limit with the set of finitely realizable
full response profiles, embeds the full response quotient into this completion,
and names the residual complement as a model-relative global realizability
defect. It then separates full finite consistency from context-restricted local
consistency and records sufficient closed-image, compactness, and shrinking-fibre
conditions under which the defect disappears.
\end{quote}

No priority claim is made for compactness, inverse limits, finite-intersection
arguments, global-section obstructions, or extension theorems themselves.

\section{Open obligations}

The next obligations are:
\begin{enumerate}
\item determine whether finite-projection closure is the most useful canonical
completion for WREH applications, or whether a uniform/metric completion is
needed alongside it;
\item study stability: approximate finite consistency under noisy responses;
\item distinguish model-class failure from protocol incompatibility;
\item formulate context families for physically incompatible measurements;
\item identify useful obstruction certificates for $\Dfin\neq\varnothing$;
\item determine when the completion defect survives admissible enlargement of
$\W$;
\item build computational examples where the defect is nontrivial but
certifiable.
\end{enumerate}

\section*{Conclusion}

WR-I asked which worlds are distinguished by available responses. WR-II asks
whether finitely compatible responses can be assembled into one admissible
global realization. The answer is not automatic. The finite-support example
shows a genuine full-finite completion defect; compactness and closed-image
criteria show when that defect disappears; context-restricted examples connect
the programme to established local-to-global obstruction theory.

The mathematical lesson is deliberately weaker than an ontological one:
\[
\text{finite consistency}
\not\Rightarrow
\text{global realization}
\]
without additional structure, while
\[
\text{compactness/closedness or shrinking fibres}
\]
can restore existence or uniqueness under explicit hypotheses.

\begin{thebibliography}{99}

\bibitem{MalachevskyWRI}
A. A. Malachevsky,
\emph{Response-Defined World Spaces: Response Equivalence,
Finite-Resolution Separation, and Identifiability of Global Realizations},
WR-I, World Realizability \& Epistemic Horizons (2026).
DOI: 10.5281/zenodo.23210258.



\bibitem{HettichKortanek1993}
R. Hettich and K. O. Kortanek,
\emph{Semi-Infinite Programming: Theory, Methods, and Applications},
SIAM Review 35(3) (1993), 380--429.
DOI: 10.1137/1035089.

\bibitem{MardesicRubin1989}
S. Marde\v{s}i\'c and L. R. Rubin,
\emph{Approximate inverse systems of compacta and covering dimension},
Pacific Journal of Mathematics 138 (1989), 129--144.
DOI: 10.2140/pjm.1989.138.129.

\bibitem{MardesicSegal1989}
S. Marde\v{s}i\'c and J. Segal,
\emph{Stability of almost commutative inverse systems of compacta},
Topology and its Applications 31 (1989), 285--299.
DOI: 10.1016/0166-8641(89)90025-4.

\bibitem{AbramskyBarbosaMansfield2017}
S. Abramsky, R. S. Barbosa, and S. Mansfield,
\emph{Contextual Fraction as a Measure of Contextuality},
Physical Review Letters 119 (2017), 050504.
DOI: 10.1103/PhysRevLett.119.050504.

\bibitem{Vorobev1962}
N. N. Vorob'ev,
\emph{Consistent Families of Measures and Their Extensions},
Theory of Probability and Its Applications 7(2) (1962), 147--163.
DOI: 10.1137/1107014.

\bibitem{AbramskyBrandenburger2011}
S. Abramsky and A. Brandenburger,
\emph{The Sheaf-Theoretic Structure of Non-Locality and Contextuality},
New Journal of Physics 13 (2011), 113036.
DOI: 10.1088/1367-2630/13/11/113036.

\bibitem{AtseriasKolaitis2023}
A. Atserias and P. G. Kolaitis,
\emph{Consistency, Acyclicity, and Positive Semirings},
in \emph{Samson Abramsky on Logic and Structure in Computer Science and Beyond},
Springer (2023), 623--668.
DOI: 10.1007/978-3-031-24117-8\_17.

\end{thebibliography}

\end{document}
